%=============================================================
% Capitolo 7 — Minimi Quadrati e Fattorizzazione QR
%=============================================================
\chapter{Minimi Quadrati e Fattorizzazione QR}
\label{ch:minimi_quadrati}

\section{Il problema dei minimi quadrati}
\label{sec:mq_problema}

Quando il numero di dati $m$ supera il numero di parametri $n$ (sistema
\emph{sovradeterminato}), in generale non esiste $x \in \R^n$ tale che $Ax = b$
con $A \in \R^{m \times n}$, $m > n$. Il problema dei \textbf{minimi quadrati}
chiede di minimizzare il \emph{residuo}:
\begin{equation}
  \min_{x \in \R^n} \|Ax - b\|_2^2 = \min_{x \in \R^n} \sum_{i=1}^m \bigl(
  (Ax)_i - b_i \bigr)^2.
  \label{eq:mq}
\end{equation}

\subsection{Applicazione: fitting di dati}

Dati $m$ punti $\{(x_i, f_i)\}_{i=1}^m$, si cerca una funzione
$y(x) = \sum_{j=1}^n \alpha_j \varphi_j(x)$ (con $\varphi_j$ funzioni base
fissate) che minimizzi:
\begin{equation}
  S = \sum_{i=1}^m \bigl(f_i - y(x_i)\bigr)^2.
  \label{eq:mq_fit}
\end{equation}
Questo è esattamente \cref{eq:mq} con $a_{ij} = \varphi_j(x_i)$.

%------------------------------------------------------------
\section{Equazioni normali}
\label{sec:equazioni_normali}
%------------------------------------------------------------

La soluzione del problema \cref{eq:mq} soddisfa le \textbf{equazioni normali}:
\begin{equation}
  A^\top A\, x = A^\top b.
  \label{eq:normali}
\end{equation}

Se $A$ ha rango pieno ($\operatorname{rank}(A) = n$), la matrice $A^\top A$ è
definita positiva e la soluzione è unica:
\begin{equation}
  x = (A^\top A)^{-1} A^\top b.
  \label{eq:sol_mq}
\end{equation}

\begin{remark}[Mal condizionamento]
La matrice $A^\top A$ ha numero di condizionamento $\kappa_2(A^\top A) =
\kappa_2(A)^2$: risolvere le equazioni normali raddoppia il numero di cifre
perse. È preferibile usare la fattorizzazione QR.
\end{remark}

%------------------------------------------------------------
\section{Fattorizzazione QR}
\label{sec:qr}
%------------------------------------------------------------

\begin{definition}[Fattorizzazione QR]
Per $A \in \R^{m \times n}$ con $m \geq n$ e $\operatorname{rank}(A) = n$, la
\emph{fattorizzazione QR} è:
\begin{equation}
  A = QR = Q \begin{pmatrix} \hat{R} \\ 0 \end{pmatrix},
  \label{eq:QR}
\end{equation}
dove $Q \in \R^{m \times m}$ è ortogonale ($Q^\top Q = I$) e
$\hat{R} \in \R^{n \times n}$ è triangolare superiore con elementi diagonali
positivi (se si richiede unicità).
\end{definition}

\subsection{Soluzione del problema dei minimi quadrati tramite QR}

Poiché la moltiplicazione per una matrice ortogonale non cambia la norma
euclidea:
\begin{align}
  \|Ax - b\|_2^2 &= \|QRx - b\|_2^2 = \|Q(Rx - Q^\top b)\|_2^2
                  = \|Rx - Q^\top b\|_2^2 \nonumber \\
  &= \left\|\begin{pmatrix}\hat{R}x - g_1 \\ -g_2\end{pmatrix}\right\|_2^2
   = \|\hat{R}x - g_1\|_2^2 + \|g_2\|_2^2,
  \label{eq:mq_qr}
\end{align}
dove $Q^\top b = \begin{pmatrix}g_1 \\ g_2\end{pmatrix}$ con
$g_1 \in \R^n$, $g_2 \in \R^{m-n}$.

Il minimo si raggiunge quando $\hat{R}x = g_1$, sistema triangolare
superiore facilmente risolubile. Il residuo minimo vale $\|g_2\|_2$.

%------------------------------------------------------------
\section{Matrici di Householder}
\label{sec:householder}
%------------------------------------------------------------

Le \textbf{riflessioni di Householder} permettono di calcolare la fattorizzazione
QR in modo numericamente stabile.

\begin{definition}[Matrice di Householder]
Data una direzione $v \in \R^m$ con $\|v\|_2 = 1$, la \emph{matrice di
Householder} è:
\begin{equation}
  H = I - 2vv^\top.
  \label{eq:householder}
\end{equation}
\end{definition}

$H$ è \emph{simmetrica} e \emph{ortogonale}: $H^\top = H$ e $H^2 = I$, quindi
$H^{-1} = H^\top = H$.

\subsection{Costruzione della fattorizzazione QR}

Dato un vettore $z \in \R^m$, si sceglie $v$ tale che $Hz = \pm\|z\|_2\,e_1$.
Applicando successivamente $n$ riflessioni:
\begin{equation}
  H_n \cdots H_2 H_1\, A = \begin{pmatrix}\hat{R} \\ 0\end{pmatrix},
  \label{eq:QR_householder}
\end{equation}
da cui $Q^\top = H_n \cdots H_1$ e $Q = H_1 \cdots H_n$ (poiché $H_k^{-1} = H_k$).

\begin{algorithm}[htbp]
  \caption{Fattorizzazione QR tramite riflessioni di Householder}
  \label{alg:qr_householder}
  \begin{algorithmic}[1]
    \Require $A \in \R^{m \times n}$ con $m \geq n$, $\operatorname{rank}(A)=n$
    \Ensure $Q \in \R^{m \times m}$ ortogonale, $\hat{R} \in \R^{n \times n}$
            triangolare superiore, $A = Q\begin{pmatrix}\hat{R}\\0\end{pmatrix}$
    \State $R \gets A$; \; $Q \gets I_m$
    \For{$k = 1$ \textbf{to} $n$}
      \State $z \gets R(k:m, k)$ \Comment{colonna $k$ dal pivot in giù}
      \State $v \gets z + \operatorname{sign}(z_1)\|z\|_2\,e_1$; \;
             $v \gets v / \|v\|_2$
      \State $H_k \gets I_{m-k+1} - 2vv^\top$
      \State $R(k:m, k:n) \gets H_k\, R(k:m, k:n)$
      \State $Q(1:m, k:m) \gets Q(1:m, k:m)\, H_k^\top$
    \EndFor
    \State \Return $Q$, $\hat{R} \gets R(1:n, 1:n)$
  \end{algorithmic}
\end{algorithm}

\subsection{Confronto: QR vs.\ equazioni normali}

\begin{table}[htbp]
  \centering
  \caption{Confronto tra i metodi per i minimi quadrati.}
  \label{tab:mq_confronto}
  \begin{tabular}{lcc}
    \toprule
    \textbf{Metodo} & \textbf{Costo} & \textbf{Stabilità} \\
    \midrule
    Equazioni normali ($A^\top A$) & $mn^2 + \frac{n^3}{3}$ & Scarsa \\
    Fattorizzazione QR (Householder) & $2mn^2 - \frac{2n^3}{3}$ & Buona \\
    SVD                              & $\mathcal{O}(mn^2)$ & Ottima (usata quando $A$ è quasi singolare) \\
    \bottomrule
  \end{tabular}
\end{table}
